\section{Finite-degree kernels and exact counters}
\label{m:sec:polynomial-kernels}
The next class requires counters as well as terminal norms. A dominant
approximation would lose cancellations needed for multivariate comparison.
For a fixed nonzero limit put
\begin{equation}\label{m:eq:J}
 J_\lambda(t)=\psi\bigl(\lambda+t-N(\lambda)\bigr),
 \qquad t\ge N(\lambda),\qquad J_0(t)=t.
\end{equation}
Thus $h_\lambda(n)=J_\lambda(N(\lambda)+n)$, and $t$ is the
whole ordinal norm.  Corollary~\ref{f:cor:closed-evaluation} gives
$J_\lambda(t)=f^\tau(t)+\tau$, where $\tau$ is the complete
replacement count.  The offset $N(\lambda)$ is retained throughout.

For integer $d\ge1$, retained numerical norm $c\ge0$, and initial control
input $t$ with $c\le f(t)$, start with suffix $\omega^d$.
At step $i\ge1$ use total norm $f^i(t)$. If the nonfinite suffix is $\eta+n$, with nonzero limit
part $\eta$ and finite $n$, replace it by $M(\eta,f^i(t)-c)$: the greatest
suffix below its \emph{limit part} of norm at most $f^i(t)-c$.
This is the limit-part convention of the principal-kernel definition; taking
the predecessor below the whole current suffix defines a different process.
Stop as soon as the suffix is finite. Let $E^d_c(t)$ be the number of
steps and
\[
 X^d_c(t)=f^{E^d_c(t)}(t).
\]
Each replacement strictly decreases an ordinal, so the process terminates.
Theorem~\ref{f:thm:maximum-predecessor} supplies the
finite algorithm and full-budget saturation. The numerical definition makes
sense even for $c$ not realized by any actual ordinal prefix.
It also permits $c>t$: the initial $t$ records the control input, rather
than the norm $c+N(\omega^d)$ of the unpadded principal and numerical mass.
The first replacement has suffix norm $f(t)-c$ by limit-part saturation,
so its total norm is exactly $f(t)$; subsequent total norms are $f^i(t)$.
Writing $Q=M(\omega^d,f(t)-c)$, the complete count is one plus the
compensated controlled count starting at $Q$.

When $t\ge c+d+1$, the numerical process agrees, from its first replacement
through the first finite exit, with the principal kernel's norm-padded startup
$\omega^d+(t-c-N(\omega^d))$. This stronger condition is required for
correspondence with the principal kernels and the genuine-prefix interpretation; it is not
imposed on the broader numerical process.

For a genuine prefix $\mu$ whose least CNF exponent is at least $d$, and
$c=N(\mu)$, splitting the controlled history gives
\begin{equation}\label{m:eq:prefix}
 J_{\mu+\omega^d}(t)=J_\mu(X^d_c(t))+E^d_c(t),
       \qquad t\ge c+d+1.
\end{equation}
The output of $X$ is still the \emph{total} norm, so the prefix mass is not
subtracted again. An admissible signed pair $(d,j)$ has $d\ge2$ and
$j\in\mathbb Z$, or $d=1$ and $j\ge0$. For each such fixed pair, put
$X^d_{c,j}(t)=f^{E^d_c(t)+j}(t)$ on a tail where the exponent is positive.
For $d=1$, the depth is $1+j$, so a negative $j$ has no positive-depth tail.
These are coherent depths: $f(X^d_{c,j})=X^d_{c,j+1}$. A negative $j$ is a
fixed number of already available predecessors under $f$; it is not a new
variable-iteration graph.

\begin{lemma}[Exact block recursion]\label{m:lem:kernel-recursion}
One has
\[
 E^1_c=1,\quad X^1_c=f,\qquad
 E^2_c(t)=f(t)/2-\lceil c/2\rceil+1.
\]
Define
\begin{equation}\label{m:eq:V}
 V_j(t)=f^{\,f(t)/2+j+1}(t).
\end{equation}
Then $X^2_{c,j}=V_{j-\lceil c/2\rceil}$ for every fixed signed $j$ on
its positive-depth tail. For $d\ge2$, the unsigned identity
$X^d_c=X^{d-1}_c\circ X^d_{c+d}$ and the count identity
\eqref{m:eq:Erec} hold on a computable tail. The signed identity
\eqref{m:eq:Xrec} holds when $d\ge3$ or $j\ge0$:
\begin{align}
 X^d_{c,j}&=X^{d-1}_{c,j}\circ X^d_{c+d},\label{m:eq:Xrec}\\
 E^d_c&=E^{d-1}_c\circ X^d_{c+d}+E^d_{c+d}.\label{m:eq:Erec}
\end{align}
\end{lemma}
\begin{proof}
For $d=1$ the first predecessor is finite. For $d=2$, the first
predecessor has $\lfloor(f(t)-c)/2\rfloor$ copies of $\omega$ and a finite
remainder; each $\omega$ copy then costs one replacement. The function $f$
is even, which gives the displayed formula.
For general $d$ write $B=f(t)$ and $B-c=dq+\rho$, $0\le\rho<d$. The first
predecessor has leading part $\omega^{d-1}q$. Its remainder is zero if
$\rho=0$, finite one if $\rho=1$, and $\omega^{\rho-1}$ if $\rho\ge2$.
This is the largest remainder \emph{below $\omega^{d-1}$} under the norm
budget, not a maximum with no ordinal bound. Removing it leaves norm $B$
with spent count one in the first two cases, or norm
$X^{\rho-1}_{B-\rho}(B)$ with count $1+E^{\rho-1}_{B-\rho}(B)$ otherwise.
The remaining $q$ degree-$(d-1)$ blocks have retained masses
$c+d(k-1)$, for $k=q,q-1,\ldots,1$. Their $X$ maps compose and their $E$
counts add. Replacing $c$ by $c+d$ preserves the initial remainder and
initial entry but stops one block earlier. This proves both recurrences,
including the coherent final depth $j$ whenever its degree-$(d-1)$
child is admissible. For $d=2$ and negative $j$, the child has no positive
depth; the exact quadratic $V$ formula applies instead. This is also a
rank-recursive exact evaluator for every displayed kernel and counter.
\end{proof}

\begin{example}[The first counter that must be retained]
For $\lambda=\omega^3$, formula \eqref{m:eq:prefix} gives
$J_\lambda(t)=X^3_0(t)+E^3_0(t)$. For $\lambda=2\omega^3$, it gives
$J_\lambda(t)=J_{\omega^3}(X^3_4(t))+E^3_4(t)$, since
$N(\omega^3)=4$. Replacing either expression by its outer $X$ alone would
change exact graph equations and could change the result of cancellation.
\end{example}

\subsection{Separation with indices allowed to grow}
We need a uniform estimate, because the block recursion creates retained
masses comparable with $f(t)$. Merely proving a tail for each fixed mass
would not license the induction.

\begin{lemma}[Uniform kernel separation]\label{m:lem:uniform}
Put $C_d=1000d^2$, $T_3=16$, and for $d\ge4$ put
$T_d=\max(T_{d-1},100C_{d-1}+100d)$. For $d\ge3$, $c,e\ge0$, signed
$j,k$, and $R\ge0$, suppose
\begin{equation}\label{m:eq:region}
 t\ge T_d,\qquad
 f(t)\ge C_d(c+e+|j|+|k|+R+1).
\end{equation}
Then the depths are defined and at least $f(t)$, and
\begin{align*}
 X^d_{c,j}(t+1)&>f^R(X^d_{e,k}(t)),\\
 c<e\quad&\Longrightarrow\quad
 X^d_{c,j}(t)>f^R(X^d_{e,k}(t)).
\end{align*}
At equal input and equal $c$, exact order is order of $j$. On a common
computed tail, equality of two rank-$d$ values occurs exactly when input,
$c$, and $j$ are identical.
\end{lemma}
\begin{proof}
We first record the elementary degree-two estimate. If
$A=f(w)/2\ge |j|+|k|+R+4$, then
\[
 V_j(w+1)>f^R(V_k(w)).
\]
Indeed $f(w+1)/2-A\ge A$, so the left iterate starts at a larger seed and
has more iterations than the right one plus $R$. At equal input, the $V$
values are ordered exactly by their indices. This argument permits indices
and required depth varying with a later induction parameter, provided the
displayed numerical bound is paid.

For $d=3$, Lemma~\ref{m:lem:residues} at moduli two and three gives
that $B=f(t)$ is even and $B\equiv1\pmod3$; write
$a=(B-1)/3$, which is odd. The three initial entries for prefix classes
$\delta=0,1,2$ are
\[
\begin{array}{c|ccc}
\delta&\text{initial }c&\text{entry norm}&\text{spent count}\\ \hline
0&B-1=3a&B&1\\
1&B=3a+1&B&1\\
2&B-2=3(a-1)+2&f(B)&2.
\end{array}
\]
Descending $c$ by three applies $V_{-\lceil(c-3)/2\rceil}$ to the preceding
entry and adds its displayed iteration length to the counter. Write
$Y^\delta_r=X^3_{3r+\delta}(t)$. At the first common entry $r=a-1$,
$Y^0_r=f(Y^1_r)$, since the two subtracted ceilings differ by one. One
further aligned block amplifies this integer input gap to any fixed
$f$-depth by the degree-two estimate. Subsequent indices are favourable:
the class-zero index is no smaller than the class-one index. Next,
$Y^2_{a-1}=f(B)$ whereas
\[
 Y^1_{a-1}=f^\ell(B),\qquad
 \ell=f(B)/2-\lceil(3(a-1)+1)/2\rceil+1.
\]
For $a$ larger than the desired depth by ten, $\ell$ exceeds that depth
plus one. Finally align class two at entry $r-1$ with class zero at entry
$r$. The initial seeds satisfy $f(B)>B$, and
$\lceil(3r-4)/2\rceil\le\lceil(3r-3)/2\rceil$, so the degree-two
estimate again applies in the favourable direction. Consequently
\[
 Y^0_r>f^L(Y^1_r),\quad
 Y^1_r>f^L(Y^2_r),\quad
 Y^2_r>f^L(Y^0_{r+1})
\]
at all required target entries, after paying the fixed depth $L$.
All smaller block entries are at least $B$, and the absolute degree-two
indices are at most $\lceil3a/2\rceil+2$ plus fixed alignment shifts;
$f(B)/2$ pays their sum and $L$. Condition \eqref{m:eq:region}, with
$C_3=9000$, gives $c,e\le B/9000$ and
$L=R+|j|+|k|+1\le B/9000$.  Each target coefficient index is at
most $B/27000$, whereas $a=(B-1)/3$.  Thus there are more than
$L+10$ blocks between the starting classes and every required target.
The displayed quadratic estimate pays the signed depth offsets as well
as $R$.  This proves the same-input statement by chaining the three
adjacent classes.

For cross-input cubic comparison, $B'=f(t+1)>2^B$ supplies more than one
extra leading block. At the smaller process's starting coefficient,
the larger process has completed a quadratic block and has seed at least
$f(B')>f(B)$, exceeding either possible smaller seed. Align the remaining
blocks, shifting the coefficient alignment by at most
$\max(\lfloor c/3\rfloor-\lfloor e/3\rfloor,0)$. Their indices are bounded
by $\lceil3a/2\rceil+2$ plus this shift, and their differences are bounded
by twice the shift plus three. The same degree-two estimate at smaller
entry at least $B$ pays each comparison. Extra remaining blocks only
increase the final norm. This proves the cubic cross-input statement.

For the induction $d\ge4$ set $L=R+|j|+|k|+1$, $B=f(t)$.
Condition \eqref{m:eq:region} leaves at least $B/(2d)>L+2$ leading
blocks. Every aligned lower-rank mass is at most $2B$, and both its pair
sum and each self-pair sum, together with $L+1$, are at most $5B$.
All smaller aligned entry norms are at least $B$. Since
$B\ge T_d\ge100C_{d-1}$,
\[
 f(B)\ge2^B>5C_{d-1}B.
\]
Thus the induction hypothesis is valid at the \emph{smaller} entry norm
even though the new masses may be comparable with $B$.

At equal $t$ and $c<e$, align the two descents at the initial leading
coefficient of the $e$ process. If no extra leading block was needed,
the $c$ process's remainder is at least as large, with smaller retained
mass. If extra blocks were needed, its last extra full block exceeds the
other initial remainder. These facts order their aligned seeds weakly.
To justify this seed comparison explicitly, maximal predecessors are
monotone in both ordinal bound and norm budget. If entry norms $v\ge w$
and masses $a,b$ satisfy
$f^i(v)-a\ge f^i(w)-b$ at every step, the first process stays above the
second at simultaneous steps, so takes at least as many steps and has at
least as large a terminal norm. Decreasing mass at equal entry is therefore
favourable. A larger entry also pays an adverse mass difference whenever
$f(v)-f(w)\ge a-b$, since successive iterate gaps do not decrease.
At aligned block $k$ the lower masses are $c+d(k-1)<e+d(k-1)$.
Equal seeds use lower-rank same-input separation; distinct seeds use
lower-rank cross-input separation. The estimates just proved pay both,
including lower self-pairs. One such block gives depth $L$, and the rest
preserve it.

At $t+1$, its first budget $B'>2^B$ gives an extra full leading block
before alignment. Its entry is at least $B'$ and retained mass at most
$B'$. The smaller initial remainder is below that block, has entry $B$,
and retained mass at most $B$. The inequality
$f(B')-f(B)>B'$ pays even an adverse retained mass; hence the larger seed
is strict. At subsequent aligned blocks the mass pair is bounded by
$4B$; lower-rank cross-input separation, including its self-pair bounds,
applies as before. This proves the cross-input statement. Monotonicity
at fixed mass extends adjacent comparison to arbitrary larger inputs.

Finally the number of remaining blocks exceeds $|j|+|k|+1$, so signed
depths leave a positive iterate and output at least $f(t)$. Taking the
maximum of the cross-pair and the two self-pair thresholds excludes every
unequal-input or unequal-mass equality. The coherent identity settles the
remaining equal-input, equal-mass case.
\end{proof}

\subsection{Counter leaders, exact ties, and residues}
For $r\ge3$ define $A_r=t$ and successively
\[
 A_{h-1}=X^h_{c+h}(A_h)\quad(h=r,r-1,\ldots,4),\qquad
 u=X^3_{c+3}(A_3),\qquad H=f(u).
\]
We call $H$ the counter's \emph{leader}; $u$ is its leader preimage.

\begin{lemma}[Counter expansion and collision]\label{m:lem:counter}
On an effective tail,
\begin{equation}\label{m:eq:counter-expand}
 E^r_c(t)=\frac H2+1-\lceil c/2\rceil
                  +\sum_{h=3}^r E^h_{c+h}(A_h),
 \qquad \frac H2<E^r_c(t)<\frac{3H}{4}.
\end{equation}
Suppose counters of ranks $r\ge s\ge3$ have equal leaders. Then they
have the same prefix index $c$. If $r=s$ their whole values are equal;
if $r>s$ the smaller-rank input is the corresponding $A_s$, and
\begin{equation}\label{m:eq:collision}
 E^r_c(t)-E^s_c(A_s)=\sum_{h=s+1}^r E^h_{c+h}(A_h).
\end{equation}
Every counter introduced on the right has rank strictly greater than $s$.
For distinct leaders, counter order agrees with leader order. Equal
leaders are ordered by counter rank, with equality at equal rank.
\end{lemma}
\begin{proof}
Iterate \eqref{m:eq:Erec} until the cubic term; its bottom quadratic count
is $f(X^3_{c+3}(A_3))/2+1-\lceil c/2\rceil$. This proves the exact
identity. A descent with initial norm $w$ and $L$ steps satisfies
$L\le f^L(w)-w$, since each iterate increases its integer argument.
Thus each residual above rank three is less than its next entry norm,
hence less than $u$. The cubic residual plus the fixed correction is also
less than $u$, since it is at most $u-A_3+1-\lceil c/2\rceil$.
It is positive on a computed tail: the number of cubic blocks tends to
infinity. Therefore
\[
 H/2<E^r_c(t)<H/2+(r-2)u<3H/4
\]
once $f(u)>4(r-2)u$ and the finite correction bound is paid.

Equality of leaders first cancels $f$, then compares the two outer cubic
kernels $X^3_{c+3}$. Lemma~\ref{m:lem:uniform} forces equality of their
prefix indices and arguments. Cancel the equal rank-four kernels next,
then rank five, up to rank $s$. If $r=s$ the original inputs are equal.
If $r>s$ the rank-$s$ input is precisely $A_s$. Subtraction of the exact
recursions gives \eqref{m:eq:collision}. Its right side is positive and tends
to infinity. For different leaders $H_1<H_2$, their being distinct values
of $f$ gives $H_2\ge2^{H_1+1}$; the half-core bounds make the order strict.
This also proves the assertions with arbitrary fixed outer shifts after
raising the computed tail.
\end{proof}

\begin{lemma}[Effective residues of kernels and counters]\label{m:lem:kernel-residues}
For each fixed $d\ge1$, $c\ge0$ and $m\ge1$, the counter $E^d_c$
has an effectively computable constant residue modulo $m$ on a computable
tail. For every admissible signed pair $(d,j)$, the same holds for
$X^d_{c,j}$ on its positive-depth tail. The $X$
residue is $r_m$. A valid counter residue is
\begin{equation}\label{m:eq:E-residue}
 E^d_c(t)\equiv
 p_d\bigl((r_{Q_d}-c)\bmod Q_d\bigr)\pmod m,
 \qquad Q_d=d!m^{d-1},
\end{equation}
where $p_d(n)$ counts partitions of $n$ into parts at most $d$.
\end{lemma}
\begin{proof}
A coherent $X$ lies in the range of $f$ with preimage at least $t$ on a
computed tail. Apply Lemma~\ref{m:lem:residues}.
For the counter we use the restricted-partition and congruent-budget
theorems of \cite{wwcounts}, with the following exact contract.  The
integer $Q_d=d!m^{d-1}$ is a period of $p_d(n)$ modulo $m$.  A
norm-padded $\omega^d$ block of initial total norm $t\ge c+d+1$,
retaining numerical mass $c$, has complete count congruent to
$p_d(S-c)$ modulo $m$ if every subsequent total budget is at least
the current total norm and congruent to one $S\ge c$ modulo $Q_d$.
The predecessor is taken below the current limit part and the segment
ends at the first finite suffix.  The budgets need not be equal.

Take
\[
 t\ge\max\{c+d+1,Q_d\},\qquad
 S=c+((r_{Q_d}-c)\bmod Q_d).
\]
The numerical kernel agrees with the genuine norm-padded block on this
tail.  Every actual budget $f^i(t)$, $i\ge1$, exceeds the preceding
total norm and is congruent to $r_{Q_d}$, by
Lemma~\ref{m:lem:residues}.  The cited transfer theorem therefore applies.
The partition period gives \eqref{m:eq:E-residue}.  Finite dynamic
programming computes $p_d(n)$ for $n<Q_d$; all smaller valid inputs
can be evaluated exactly.  This proves an explicit residue tail without
using comparison or an existential decision procedure.
\end{proof}

\section{Closed existential word arithmetic}
\label{m:sec:word-arithmetic}
Fix a finite maximum degree $D\ge3$. A \emph{word} in one raw variable
is a finite composition of signed constant shifts and the functions
\[
 f,\quad V_j,\quad X^r_{c,j},\quad E^r_c\quad(3\le r\le D),
\]
with fixed integer indices satisfying the preceding conventions. A
word-linear row is an integer linear combination of finitely many such
words, evaluated in possibly different raw variables, plus a constant.
We allow arbitrary Boolean combinations of row comparisons with zero and
fixed-modulus congruences. Each raw variable has an explicit individual
lower bound ensuring that all its words are defined.
These bounds concern every displayed occurrence, including an occurrence
whose coefficient is zero.  The domain of the whole query is the
conjunction $\bigwedge_i x_i\ge b_i$ of those individual bounds;
it is not changed by a Boolean negation inside the matrix.

Every word has a computable domain tail and a lower offset $M$ with
$W(n)\ge n-M$. This follows structurally: shifts add to the offset;
$f$, $V$, and coherent $X$ dominate their input on their stated tails;
counters do so as well. For the latter, a sufficiently large degree-three
block contains two quadratic blocks, and a quadratic block entered at
least at $f(t)$ already spends more than $f(t)$ steps. Higher-degree blocks
contain one of these. These structural bounds also pay every threshold at
an interior argument of a composition. They are not obtained by treating
the interior argument as a raw variable.

\begin{lemma}[Ordered maximum-variable guards]\label{m:lem:guards}
For two words $A(x),B(y)$ and the guard $x\ge y$, an effective common
tail permits their three-way comparison to be rewritten as a finite Boolean
combination in which $x$ occurs only linearly and words are applied only
to the remaining variable $y$. For two words at the \emph{same} raw
variable the corresponding tail comparison is effectively decided.
All thresholds and newly generated words are computed before the tail is
used.
\end{lemma}
\begin{proof}
Peel fixed outer shifts. Unequal range values have gaps eventually larger
than a fixed shift. Lemma~\ref{m:lem:counter} gives the same fact for two
counter values, including equal-leader unequal-rank values whose exact
residual difference tends to infinity. For a counter with preimage $u$
and range value $f(v)$, its half-core bounds give
\[
 E^r_c<f(v)\ \Longleftrightarrow\ u\le v,
 \qquad E^r_c>f(v)\ \Longleftrightarrow\ u>v;
\]
the gaps again tend to infinity. Thus fixed shifts are kept for exact-core
ties and otherwise cannot close the gap.

For $f/f$ strip the two $f$ nodes. For $f$ against a coherent range
kernel, move one fixed depth to its index and strip $f$; for example
$V_j=f\circ V_{j-1}$. Two $V$ heads compare lexicographically by argument,
then by index, using the elementary estimate in the proof of
Lemma~\ref{m:lem:uniform}. Two equal-rank $X$ heads compare by argument,
then by $(-c,j)$, using that lemma and retaining the exact equal-input
branch.

For $X^r$ against $X^s$ with $r>s\ge3$, expand only the higher-rank head
by \eqref{m:eq:Xrec} until its outer rank is $s$, and immediately compare
and strip that head against the opposite rank-$s$ head. If $s=2$,
the exposed quadratic head is $V_{j-\lceil c/2\rceil}$ and the same
procedure uses its index comparison. No inverse counter is introduced.
For counter/range comparison use the displayed preimage rule. For
counter/counter comparison use their leader preimages followed by the
rank and exact-tie rule of Lemma~\ref{m:lem:counter}.

Here is a termination measure, needed because mixed-rank expansion can
increase the written expression. First count all $E$ nodes in the two
words. Replacing a counter by its leader-preimage composition removes
its outer $E$ node and retains its argument once. For the range part,
assign each rank-$s$ occurrence weight $\omega^{D-s}$, where $f$ has
rank one and $V$ rank two, and naturally sum these weights. A mixed
comparison deletes the opposite rank-$s$ occurrence and introduces only
finitely many occurrences of ranks strictly above $s$, retaining the
original higher-rank occurrence. Its ordinal weight decreases. Equal-rank
stripping and $f$ cancellation decrease weight as well. The lexicographic
pair (number of $E$ nodes, range weight) is therefore well-founded along
every recursive child. There are finitely many branches at each step.

Ultimately the $x$ side is a shifted leaf, so its comparison with a word
at $y$ is linear in $x$; or the $y$ side is a shifted leaf and the $x$
side is nonlinear. In the latter case the nonlinear value is at least
$2^{x-M}$ for a computed $M$, so eventually exceeds $y$ plus any fixed
shift, since $y\le x$. When both words have the same raw variable,
continue stripping until leaves or leaf/nonlinear growth decides the
comparison; do not create a purported linear atom $x<W(x)$.
Take the finite maximum of all local bounds using the structural offsets
of interior arguments. The finite recursion tree, rather than a guessed
fixed envelope of indices, computes that maximum.
\end{proof}

The guards compare individual word values. A linear row can still cancel
its largest values, so it requires a different induction.

\begin{lemma}[Ranked cancellation]\label{m:lem:row}
For a finite collection of word-linear rows and a chosen maximal raw
variable $x$, one can compute a tail and finite guarded cases in which
each row either has a known strict sign or is exactly linear in $x$.
The guards can themselves be made linear in $x$ by
Lemma~\ref{m:lem:guards}. No correlation between word values is discarded.
\end{lemma}
\begin{proof}
Absorb fixed outer shifts into the constant and keep raw-linear terms
separate. The other top-level atoms are ordinary range values (rank zero
for this row argument) and counters of ranks $3,\ldots,D$. Partition their
ordinary values and counter leaders into ordered equality classes, retaining
every equality and order guard. Merge equal-rank counters in a class as
\emph{whole counters}, justified by Lemma~\ref{m:lem:counter}. Prefer a
representative actually involving $x$ when one remains in the current
class. A class now has ordinary coefficient $a$ and at most one coefficient
$b_r$ for each counter rank. Its leading coefficient is
\[
 a+\tfrac12\sum_r b_r.
\]
If it vanishes, pay its residuals before deciding any sign. When $a\ne0$,
expand the counters by \eqref{m:eq:counter-expand}; the common half-leader
cancels the ordinary atom, leaving only residual counters and constants.
When $a=0$, let $s$ be the least surviving counter rank. Rewrite each
higher counter by \eqref{m:eq:collision} as the common rank-$s$ counter plus
higher-rank residuals. Since $\sum b_r=0$, the common term cancels. Remove
an entirely zero class. Repartition after every rewrite, since its newly
exposed leaders may coincide with leaders elsewhere in the row.

Count \emph{top-level} row atoms, assigning the ordinary rank-zero atom
weight $\omega^D$ and a rank-$r$ counter weight $\omega^{D-r}$.
Natural sum is a well-founded measure. The first rewrite loses a rank-zero
atom and adds only positive-rank atoms. The second loses the least rank
$s$ and adds only ranks strictly greater than $s$. Existing higher-rank
atoms can proliferate, but cannot offset that ordinal decrease. Nested
nodes inside their arguments are not included in this measure; their
comparisons terminate by the separate guard measure. Every branching row
rewrite therefore strictly decreases weight. The tree is finitely
branching and has no infinite branch, so it is finite. Enumerating it to
completion computes all generated words, coefficients, prefix indices,
and guard thresholds before any finite-prefix case is split off.

After the rewrites every surviving leading coefficient has absolute value
at least $1/2$. Let $P=f(u)$ be the greatest surviving leader, $S$ the
current absolute nonlinear coefficient mass, and $A,B$ bounds for the
absolute raw-linear coefficients and constant. The exact counter
corrections have total size at most $DSu$. Every smaller leader $Q$
satisfies $Q<\log_2 P$, so its total contribution is at most
$S\log_2P$. For $u\ge8DS+16$ each of these two errors is less than $P/16$.

If a current nonlinear $x$ occurrence survives, the greatest leader is
at least its leader. Structural offsets therefore give $u\ge x-M$ for a
computed $M$. Require $x\ge2M+16$ and $x\ge8(A+B+1)$. The elementary
inequality $2^{x-M}\ge x^2$ on this range, proved at the initial value
and then by doubling, gives $Ax+B<P/8$. Total error is less than $P/4$,
while the leading magnitude is at least $P/2$. Hence its coefficient
strictly determines the sign. If no current nonlinear $x$ occurrence
remains, keep the exact row, which is linear in $x$.

Growth provenance refers only to current occurrences. A deleted leader
passes no $x$ label to a residual which does not actually contain $x$.
This prevents an invalid sign decision after cancellation has removed all
nonlinear dependence on $x$. Take maxima of these explicit bounds over the
finite row and guard trees. Ordered partitions retain the original shared
arguments throughout, even when an inconsistent branch is later discarded.
\end{proof}

\begin{theorem}[Closed existential word arithmetic]\label{m:thm:existential}
The existential closure of every Boolean combination of word-linear row
comparisons and fixed-modulus congruences in the finite-degree language
above is decidable, on its explicit word domains.
\end{theorem}
\begin{proof}
First remove a variable which occurs only linearly, if any, by ordinary
Presburger quantifier elimination \cite{cooper}, including its explicit
lower bound in the matrix. Words in other variables are parameters
in this step; the eliminated variable occurs in none of them.  The
Presburger equivalence is valid for every assignment of its parameter
symbols, hence for their actual correlated word values.  Its output
uses only integer-linear rows and fixed-modulus congruences in those
same parameters.  Substituting the words back therefore stays in the
stated syntax and does not apply any word to a new sum. Repeat or
evaluate the closed matrix if no variables remain.

Otherwise branch on a maximal raw variable $x$, retaining $x\ge x_i$.
Compute the complete finite guard and row trees of
Lemmas~\ref{m:lem:guards}--\ref{m:lem:row}. Compute residue thresholds for each
congruence as well: Lemma~\ref{m:lem:kernel-residues}, composed with interior
lower-offset bounds, makes every nonlinear word's residue constant.
Collect shifts in its fixed residue. On that tail, congruences are thus
congruences in raw variables only. Take a common $N$ paying all original
and generated domains, sign estimates, guards, and residue bounds.

Before executing any tail transformation, split the \emph{original current
matrix} into the all-raw-inputs-$\ge N$ branch and finitely many branches
fixing some raw $x_i=c<N$ above its original lower bound. In a finite
branch substitute $c$ into its original words and evaluate them exactly.
Do not evaluate a newly introduced coherent negative-depth word below the
tail on which it was introduced. Exact kernel evaluation is finite by the
block recursion. Each such branch removes a quantified variable.

On the all-large branch, use the guarded row and congruence transformations.
Everything now contains $x$ only linearly. Presburger elimination of
$\exists x$, treating words in the remaining variables as parameters,
again produces word-linear rows and congruences with one fewer variable.
Preserve all Boolean connectives and every equality/order guard. No word
has been applied to a sum, and no auxiliary derived value has been made
an independent witness.

The output retains the lower bounds of the remaining variables.  Each
new word introduced in the finite normalization tree has already had
its genuine domain paid by the chosen common $N$.  Thus the recursive
call has the same explicit individual-domain contract as the input.
Any new congruence modulus created by Presburger elimination is dealt
with in that smaller-variable call, not by restarting the present
dimension in search of a stable cutoff.

The induction measure for the whole algorithm is the number of raw
existential variables. It strictly decreases in every finite-prefix,
pure-linear, and all-large branch. The subsidiary guard and row trees are
finite by their own measures. Consequently this is a terminating effective
procedure. No feasible complexity estimate follows: finite exception values
and arithmetic thresholds can be extremely large.
\end{proof}

This theorem is deliberately about \emph{closed existential} sentences.
In a formula with a free parameter, that parameter could be the largest raw
argument and could not be removed by the stated induction. We have proved
neither free-parameter quantifier elimination nor the full first-order
theory of arbitrary polynomial-collapse functions.

\subsection{Exact polynomial graphs}
\begin{theorem}[Polynomial collapse graphs]\label{m:thm:graphs}
For each fixed nonzero polynomial limit $\lambda<\omega^\omega$, the graph
$y=h_\lambda(n)$ is effectively expressible as a finite disjunction of
word-linear equations in $n,y$, with explicit domains and exact finite
exceptions. One finite degree bound suffices for each finite collection of
such graphs. Arbitrary fixed nesting and addition of these functions admit
a total functional existential graph translation.
\end{theorem}
\begin{proof}
Remove one unit at the least positive degree of $\lambda$, writing
$\lambda=\mu+\omega^d$, and put $c=N(\mu)$. There are finitely many
monomial units. Starting with $2J_0(t)=2t$, compile \eqref{m:eq:prefix}
recursively: substitute the unary word $X^d_c(t)$ for the input in the
already compiled expression for $2J_\mu$, and add $2E^d_c(t)$. For $d=1$
use $X=f$, $2E=2$; for $d=2$ use
\[
 X=V_{-\lceil c/2\rceil},\qquad
 2E=f(t)-2\lceil c/2\rceil+2;
\]
for $d\ge3$ retain the $X$ and $E$ heads. This substitution composes unary
words only. It never applies a function to a sum of compiled values.

Finally replace $t$ by $n+N(\lambda)$. On a computed tail this yields
$2y=P_\lambda(n)$ with an integer word-linear $P_\lambda$. The finitely
many excluded $n\ge0$ are evaluated by the total closed predecessor
algorithm and included as disjuncts $n=k$, $y=h_\lambda(k)$. The graph is
therefore total and functional on the original natural domain.
To flatten a nested expression, introduce a fresh natural variable for
each arithmetic node: addition and fixed multiples have linear equations;
a collapse node has the graph just constructed, with its original argument
variable. Preserve all repeated original variable names and quantify every
new node variable. This creates no independence assumption: the graph
equations enforce precisely the original correlations.
\end{proof}

\begin{lemma}[From finite graph exceptions to word domains]
\label{m:lem:graph-domains}
Let finitely many total functional graph clauses be obtained from
Theorem~\ref{m:thm:graphs}, together with linear arithmetic graph
equations and a Boolean combination of linear comparisons between graph
wires.  Their closed
existential conjunction is decidable by
Theorem~\ref{m:thm:existential}, with original positive inputs and
fresh natural graph variables.
\end{lemma}
\begin{proof}
Every collapse graph is a finite disjunction of its tail equation,
guarded by a computed input lower bound, and finitely many exact
small-input point clauses.  Distribute only these finitely many
positive graph disjunctions.  In a resulting case, a small-input
clause fixes its input and output to known integers.  Substitute such
equalities into the case.  If a retained tail guard on the same input
fails at the substituted integer, discard the inconsistent case;
otherwise any word in that fixed input is defined and can be evaluated
exactly.  Repeat the finite substitutions if desired.

All remaining word occurrences come from tail clauses with their
explicit bounds.  Take, for each shared raw variable, the maximum of
the bounds required by its remaining words, together with its original
natural or positive domain.  The accompanying guards assert exactly
these bounds.  The case is now a closed existential Boolean combination
of word-linear rows and congruences on the individual word domains, so
Theorem~\ref{m:thm:existential} applies.  The disjunction is finite.

This argument does not replace a correlated node value by a free
parameter.  All linear graph equations and original names remain in
each case.  It also does not negate a graph: the output test may contain
negations, but the functional graph constraints remain positively
conjoined.  Their uniqueness on the original domains gives the intended
source semantics.
\end{proof}

\begin{corollary}[Polynomial-head Boolean comparison]
\label{m:cor:polynomial-booleans}
For positive expressions in any finite number of raw inputs, using only
fixed polynomial limit heads below $\omega^\omega$, the existential
and universal closures of every finite Boolean combination of equality
and order atoms are decidable.  Arbitrary fixed nesting and addition
are permitted.
\end{corollary}
\begin{proof}
Flatten all the finitely many input circuits together, assigning one
fresh natural wire to each node and one shared input to each repeated
raw name.  Conjoin their total functional graphs positively and append
the Boolean output test.  The graph constraints have a unique extension
at every original positive valuation, so existence of this full
conjunction is equivalent to the original existential query.
Lemma~\ref{m:lem:graph-domains} decides it.  For a universal query,
negate only the Boolean output test and decide the resulting existential
counterexample; the graph conjunction remains positive.  All fresh
wires and original inputs belong to the same existential block.
\end{proof}

